\documentclass[11pt]{article}

\usepackage[margin=1in]{geometry}
\usepackage{booktabs}
\usepackage{hyperref}

\title{11-768 HW2: Evaluating a Data-Visualization Agent}
\author{Your Name}
\date{}

\begin{document}
\maketitle

\section{Part 1: Examine Seed Agent Runs}

\subsection{Comparison of Ground-Truth and Validator Labels}
Summarize how the baseline validator's predictions compare with the provided
human labels on the seed runs. Report all four per-family MCC scores and macro-MCC, and describe
where the disagreements concentrate.

\subsection{Validator Error Type 1}
Describe the first recurring type of validator error. Cite the run IDs of at
least two runs that exhibit it, and explain why the baseline validator failed
on them. Ground the analysis in specific evidence from the task, trajectory,
input data, or figure.

\subsection{Validator Error Type 2}
Describe the second recurring type of validator error. Cite the run IDs of at
least two runs that exhibit it, and explain why the baseline validator failed
on them. Ground the analysis in specific evidence from the task, trajectory,
input data, or figure.

\subsection{Proposed Validator Changes}
For each error type above, propose a targeted change to the validator that
might catch that class of error.

\section{Part 2: Improve Your Validator}

\subsection{Implemented Changes}
Describe the two changes implemented from Part 1. Distinguish deterministic
checks from model-based checks, where applicable, and explain how each change
uses the task, trajectory, input data, or rendered figure.

\subsection{Seed-Set Performance vs.\ Baseline}
Report the modified validator's performance on the seed runs alongside the
baseline using all four per-family MCC scores and macro-MCC.

\subsection{Qualitative Analysis of Targeted Error Types}
For each of the two error types surfaced in Part 1, discuss whether the
corresponding change actually mitigated it, with reference to specific runs.
Also describe any regressions.

\section{Part 3: Design New Evaluation Tasks}

\subsection{Anticipated Generalization Weaknesses}
Identify two classes of visualization task that your current validator may
have trouble generalizing to, and explain why.

\subsection{Task Design}
Describe the five visualization tasks, including at least two from each
weakness class. For each task, describe its inputs and data source, expected
behavior, and why it is a useful stress test for the validator.

\subsection{Agent Runs and Labeling}
Describe the agents run on the new tasks, the runs generated, and how you
assigned human labels using the four fixed error families. Explain any
ambiguous labeling decisions.

\subsection{Validator Performance on New Tasks}
Report all four per-family MCC scores and macro-MCC on the new agent runs and qualitatively
analyze important successes, failures, and uncertain labels.

\subsection{Harbor Packaging}
State that every authored task scores 1.0 against \texttt{-a oracle}; you may
also report the recommended \texttt{-a nop} runs. Describe both plausible-but-wrong
mutant solutions and report each reward. For the mutant scoring 0, explain what
it gets wrong and name the assertion that rejects it. For the mutant scoring
1.0, explain what it gets wrong and why the verifier accepts it. Then, for each
of the four error families, say whether a deterministic Harbor verifier for your
task could decide it, and why or why not. Finally: your validator is graded on
private runs for tasks you never authored. Explain why a per-task verifier
cannot serve as that grading mechanism.

\section{Part 4: Iterate on Your Design}

\subsection{Additional Improvement Motivated by Part 3}
Describe at least one change to the validator motivated by the Part 3 results.
Identify the observed failure that motivated it and explain why you expected
the change to help.

\subsection{Final Evaluation on Seed and New Tasks}
Evaluate the final validator on both the seed runs and self-authored runs.
Report all four per-family MCC scores and macro-MCC before and after the change on each subset,
including any regressions.

\subsection{Qualitative Analysis of Self-Authored Tasks}
Discuss whether the change improved performance on your self-authored tasks,
with reference to specific runs.

\subsection{Further Changes}
Report any additional validator changes or tasks beyond the required ones.
For each, explain its motivation and how successful you judged it to be,
qualitatively and, for quantitative comparisons, using all four per-family MCC
scores and macro-MCC.

\end{document}
